%%==============================================================================
%% sections/05_experiments.tex
%%==============================================================================
\section{Experiments}
\label{sec:experiments}

\subsection{Setup}
\paragraph{Training compute} All models are trained on 64$\times$NVIDIA H800
GPUs. The 7B model consumes \SI{3.2e22}{FLOPs}, roughly $3.1\%$ of the
Llama-3-70B budget.

\paragraph{Benchmarks} We evaluate on six standard reasoning benchmarks:
GSM8K, MATH, HumanEval, MBPP, ARC-Challenge, and BBH.

\subsection{Main Results}
\begin{table}[t]
    \centering
    \small\sffamily
    \begin{tabularx}{\linewidth}{lXcccc}
        \toprule
        \textbf{Model} & \textbf{Params} & \textbf{GSM8K} & \textbf{MATH} & \textbf{HumanEval} & \textbf{Avg.} \\
        \midrule
        Llama-3-8B         & 8\,B  & 76.3 & 46.4 & 67.1 & 63.3 \\
        Mistral-7B         & 7\,B  & 68.5 & 41.2 & 62.8 & 57.5 \\
        Qwen-2-7B          & 7\,B  & 79.2 & 50.1 & 68.7 & 66.0 \\
        \midrule
        \textbf{BJTU-LLM-7B (ours)} & \textbf{7\,B} & \best{82.4}\,\up{6.1} & \best{54.7}\,\up{8.3} & \best{71.3}\,\up{4.2} & \best{69.5} \\
        \bottomrule
    \end{tabularx}
    \caption{Main results on reasoning benchmarks (pass@1, greedy decoding).
    Gains vs.\ Llama-3-8B shown in blue. Best per column in \textbf{bold}.}
    \label{tab:main}
\end{table}

\begin{keypoint}
\textbf{Key finding.} BJTU-LLM-7B is the strongest 7B model across all three
reasoning benchmarks, and closes 78\% of the gap to Llama-3-70B.
\end{keypoint}

\subsection{Ablation Study}
\begin{table}[t]
    \centering
    \small\sffamily
    \begin{tabular}{lcc}
        \toprule
        \textbf{Configuration} & \textbf{GSM8K} & \textbf{$\Delta$} \\
        \midrule
        Full BJTU-LLM-7B                       & 82.4 & --- \\
        \quad w/o synthetic-data curriculum     & 74.1 & \dn{8.3} \\
        \quad w/o sparse attention              & 81.8 & \dn{0.6} \\
        \quad w/o process rewards (outcome only) & 78.6 & \dn{3.8} \\
        \bottomrule
    \end{tabular}
    \caption{Ablation on GSM8K. Each component is removed independently.}
    \label{tab:ablation}
\end{table}

\subsection{Qualitative Analysis}
Figure~\ref{fig:qual} shows a representative example of BJTU-LLM-7B solving a
multi-step math problem via chain-of-thought reasoning.

\begin{prompt}[Prompt]
A train leaves Beijing at 8:00 traveling south at 250\,km/h. Another train
leaves Shanghai at 9:00 traveling north at 300\,km/h. If the two cities are
1318\,km apart, at what time do the trains meet?
\end{prompt}

\begin{response}[BJTU-LLM-7B]
Step 1: At 9:00, the first train has traveled $250 \times 1 = 250$ km.\\
Step 2: Remaining distance: $1318 - 250 = 1068$ km.\\
Step 3: Closing speed: $250 + 300 = 550$ km/h.\\
Step 4: Time to close: $1068 / 550 \approx 1.94$ h $=$ 1 h 56 min.\\
Answer: The trains meet at approximately \textbf{10:56}.
\end{response}

\begin{figure}[t]
    \centering
    \fbox{\parbox{0.9\linewidth}{\centering\footnotesize
    (Placeholder: attention heatmap / accuracy-vs-compute curve.
     Drop your PDF/PNG into \texttt{figures/} and use
     \texttt{\textbackslash includegraphics\{filename\}}.)}}
    \caption{Qualitative visualization of reasoning behavior.}
    \label{fig:qual}
\end{figure}
